\section{Report Structure}
\label{sec:intro_report_structure}

In the following chapters, we cover the following:

\paragraph{Background:}

\noindent Chapter~\ref{cha:background} provides the theoretical context for the \texttt{DetectorNX} project. It surveys \glspl{cnn} for object detection, introduces \glspl{snn} via the \gls{lif} neuron and surrogate gradients, and reviews event-based vision alongside the \gls{etram} dataset. It then characterises neuromorphic energy economics, concluding with a synthesis of related work that identifies the critical gap in high-definition spiking object detection.

\paragraph{Methodology:}

\noindent Chapter~\ref{cha:methodology} presents the full technical implementation.
It details the system requirements, the shared multi-box detection head architecture, the data preprocessing pipeline, the design of \texttt{DetectorNX-G3-SNN} and \texttt{DetectorNX-G3-CNN}, the architectural evolution that led to the final configurations, and the complete experimental setup including training hyperparameters, hardware environment, and evaluation protocol.
Results and their analysis are presented within this chapter as an integrated discussion of methodology and outcome.

\paragraph{Conclusions:}

\noindent Chapter~\ref{cha:summary_conclusions} concludes the report.
It revisits each success criterion defined in Section~\ref{sec:intro_success_criteria}, summarises the principal findings across the accuracy, latency, and energy dimensions, reflects critically on the limitations of the experimental design, and proposes concrete directions for future work including deployment on physical neuromorphic hardware and exploration of alternative spiking encodings.
